\section{Results}
\subsection{Inference: odds of robbery by time}
Figure~\ref{fig:forest} shows the odds ratios. After 2019 the odds that a reported property-type incident is a robbery are about half what they were before: the post-2019 odds ratio is \resG{post2019.or} (robust \pct{} interval \resG{post2019.robust_or_lo} to \resG{post2019.robust_or_hi}). The composition check matters here. Fitting robbery against each other category separately gives \resG{composition.Auto Theft.post2019.or} against auto theft, \resG{composition.Break and Enter.post2019.or} against break and enter and \resG{composition.Theft Over.post2019.or} against theft over, so the pooled ratio mixes a fall in robbery with a large rise in auto theft, the comparison class. Time of day matters most. Against noon, the odds of robbery are highest at 15:00 (\resG{hour_or_vs_12.15.or}, interval \resG{hour_or_vs_12.15.robust_or_lo} to \resG{hour_or_vs_12.15.robust_or_hi}) and lowest at 08:00 (\resG{hour_or_vs_12.8.or}, \resG{hour_or_vs_12.8.robust_or_lo} to \resG{hour_or_vs_12.8.robust_or_hi}). After dark the odds are higher, with ratio \resG{dark.or} (\resG{dark.robust_or_lo} to \resG{dark.robust_or_hi}). Month and weekday terms are small: the largest month effect is December against January (\resG{month_or_vs_jan.12.or}). Of the \resI{time_terms} time terms, \resI{time_terms_clear} have robust intervals that exclude one. Even so, the model separates the classes only weakly: its cross-validated AUC is \resG{cv.auc_mean} (standard deviation \resG{cv.auc_sd} over folds).

\begin{figure}[t]
\centering
\includegraphics[width=\linewidth]{build/figures/fig2.pdf}
\caption{Odds ratios of robbery versus other major property crime, with cluster-robust \pct{} intervals on a log scale. Blue marks the post-2019 period, darkness and 15:00. Many hours and a few months have intervals that exclude one (the dashed line); most weekdays do not.}\label{fig:forest}
\end{figure}

\subsection{Prediction: what each ingredient adds}
Table~\ref{tab:metrics} and Figure~\ref{fig:ladder} give the ladder. Time alone is a weak predictor. Under Design~A it reaches a PR-AUC of \resA{variants.L1.pr_auc.estimate} \ci{\resA{variants.L1.pr_auc.lo}}{\resA{variants.L1.pr_auc.hi}} against a baseline of \resA{prevalence}, and a ROC-AUC of \resA{variants.L1.roc_auc.estimate}. Adding place is the big step: PR-AUC rises by \resA{deltas.L2_minus_L1.pr_auc.estimate} \ci{\resA{deltas.L2_minus_L1.pr_auc.lo}}{\resA{deltas.L2_minus_L1.pr_auc.hi}} to \resA{variants.L2.pr_auc.estimate}. Most of that step is premises and location type: removing them from L2 costs \resA{deltas.L2_minus_L2-nopremises.pr_auc.estimate} of PR-AUC \ci{\resA{deltas.L2_minus_L2-nopremises.pr_auc.lo}}{\resA{deltas.L2_minus_L2-nopremises.pr_auc.hi}}. The census block adds little once place is known: \resA{deltas.L3_minus_L2.pr_auc.estimate} of PR-AUC \ci{\resA{deltas.L3_minus_L2.pr_auc.lo}}{\resA{deltas.L3_minus_L2.pr_auc.hi}} and \resA{deltas.L3_minus_L2.roc_auc.estimate} of ROC-AUC \ci{\resA{deltas.L3_minus_L2.roc_auc.lo}}{\resA{deltas.L3_minus_L2.roc_auc.hi}}. Non-linear models add more (random forest over L3: \resA{deltas.L4_minus_L3.pr_auc.estimate} of PR-AUC), but gradient boosting adds nothing measurable over the forest: the paired difference in PR-AUC has interval \resA{deltas.L5_minus_L4.pr_auc.lo} to \resA{deltas.L5_minus_L4.pr_auc.hi}, which includes zero. Without premises and location type even the best model falls to a PR-AUC of \resA{variants.L5-nopremises.pr_auc.estimate}.

\begin{table}[t]
\centering\footnotesize\tabsetup
\caption{Ladder metrics with \pct{} cluster-bootstrap intervals (ROC-AUC, PR-AUC, F1) and point estimates (precision, recall at the tuned threshold). F1, precision and recall use a threshold tuned on inner out-of-fold predictions and are slightly optimistic. Prevalence differs between designs, so PR-AUC is compared only within a design.}\label{tab:metrics}
% Generated by tools/gen_paper.py from results/. Do not edit by hand.
\begin{tabular}{@{}l@{\hskip 4pt}ccccc@{}}
\toprule
Model & ROC-AUC & PR-AUC & F1 & Precision & Recall \\
\midrule
\multicolumn{6}{@{}l}{\textit{Neighbourhood-grouped 5-fold cross-validation}} \\
L1 Time only & 0.588 [0.579, 0.597] & 0.197 [0.179, 0.214] & 0.300 [0.278, 0.320] & 0.183 [0.166, 0.198] & 0.832 [0.823, 0.843] \\
L2 + place & 0.836 [0.824, 0.847] & 0.531 [0.504, 0.555] & 0.546 [0.522, 0.565] & 0.483 [0.448, 0.510] & 0.628 [0.607, 0.647] \\
L3 + census & 0.839 [0.830, 0.849] & 0.543 [0.518, 0.566] & 0.549 [0.530, 0.568] & 0.490 [0.467, 0.512] & 0.625 [0.604, 0.646] \\
L4 Random forest & 0.858 [0.846, 0.869] & 0.593 [0.561, 0.620] & 0.562 [0.538, 0.583] & 0.520 [0.486, 0.549] & 0.611 [0.590, 0.631] \\
L5 XGBoost & 0.860 [0.851, 0.869] & 0.593 [0.564, 0.618] & 0.564 [0.545, 0.582] & 0.516 [0.493, 0.538] & 0.622 [0.599, 0.645] \\
L2 + place, no premises & 0.599 [0.576, 0.620] & 0.202 [0.174, 0.232] & 0.306 [0.282, 0.326] & 0.194 [0.175, 0.210] & 0.721 [0.681, 0.759] \\
L5 XGBoost, no premises & 0.643 [0.626, 0.658] & 0.242 [0.215, 0.273] & 0.330 [0.303, 0.352] & 0.220 [0.200, 0.237] & 0.655 [0.603, 0.702] \\
\midrule
\multicolumn{6}{@{}l}{\textit{Temporal holdout (train to 2021, test 2022 to 2023)}} \\
L1 Time only & 0.586 [0.574, 0.598] & 0.141 [0.126, 0.157] & 0.215 [0.195, 0.235] & 0.125 [0.112, 0.138] & 0.775 [0.760, 0.790] \\
L2 + place & 0.808 [0.798, 0.819] & 0.362 [0.336, 0.388] & 0.385 [0.362, 0.407] & 0.313 [0.290, 0.337] & 0.499 [0.472, 0.525] \\
L3 + census & 0.813 [0.802, 0.823] & 0.377 [0.349, 0.404] & 0.390 [0.366, 0.412] & 0.318 [0.295, 0.342] & 0.504 [0.478, 0.530] \\
L4 Random forest & 0.840 [0.830, 0.849] & 0.440 [0.411, 0.468] & 0.438 [0.415, 0.459] & 0.371 [0.349, 0.393] & 0.533 [0.499, 0.563] \\
L5 XGBoost & 0.841 [0.830, 0.849] & 0.435 [0.406, 0.463] & 0.437 [0.415, 0.458] & 0.352 [0.331, 0.373] & 0.575 [0.543, 0.606] \\
L2 + place, no premises & 0.610 [0.593, 0.627] & 0.158 [0.136, 0.184] & 0.233 [0.211, 0.253] & 0.141 [0.127, 0.155] & 0.662 [0.612, 0.704] \\
L5 XGBoost, no premises & 0.665 [0.642, 0.684] & 0.202 [0.171, 0.234] & 0.264 [0.246, 0.284] & 0.168 [0.156, 0.181] & 0.616 [0.552, 0.678] \\
\bottomrule
\end{tabular}

\end{table}

The temporal holdout tells the same story at lower levels, as expected when the robbery share has fallen to \resB{prevalence}. Time alone gives a PR-AUC of \resB{variants.L1.pr_auc.estimate} \ci{\resB{variants.L1.pr_auc.lo}}{\resB{variants.L1.pr_auc.hi}}, close to that baseline, and L5 reaches \resB{variants.L5.pr_auc.estimate} (ROC-AUC \resB{variants.L5.roc_auc.estimate} against \resA{variants.L5.roc_auc.estimate} under Design A).

\begin{figure}[t]
\centering
\includegraphics[width=\linewidth]{build/figures/fig3.pdf}
\caption{PR-AUC by rung (L1 time only; L2 adds place; L3 adds census; L4 random forest; L5 XGBoost) under neighbourhood-grouped cross-validation (left) and the temporal holdout (right), with \pct{} cluster-bootstrap intervals. Diamonds drop premises and location type. Dotted lines show the share of robberies, the PR-AUC of a constant score.}\label{fig:ladder}
\end{figure}

\subsection{Inference versus prediction}
Figure~\ref{fig:contrast} puts the two analyses side by side. In the regression, \resI{inference.hour.clear} of the \resI{inference.hour.terms} hour terms and every darkness and post-2019 term are statistically clear. In the prediction ladder, the time block matters little: shuffling it costs the L5 model \resA{variants.L5.block_permutation_importance_pr_auc_drop.time} of PR-AUC under Design~A, against \resA{variants.L5.block_permutation_importance_pr_auc_drop.premises} for premises and location type, and only \resA{variants.L5.block_permutation_importance_pr_auc_drop.census} for census and \resA{variants.L5.block_permutation_importance_pr_auc_drop.geo} for geography. Both findings are true at once. With \resG{n} rows, small associations get narrow intervals; a term can be clear and still move the odds too little to sort individual incidents.

\begin{figure}[t]
\centering
\includegraphics[width=\linewidth]{build/figures/fig4.pdf}
\caption{Left: for each time feature, the number of terms in the regression whose robust interval excludes one (labels give the count of clear terms out of all terms). Right: PR-AUC lost by the L5 model on held-out neighbourhoods when a block of features is shuffled; ``Time'' covers hour, month, weekday and darkness.}\label{fig:contrast}
\end{figure}
